	\subsubsection{Perímetro y Área}

	\ejemplosboxed[6]{Encuentra el \textbf{perímetro} y el \textbf{área} de las siguientes figuras:

		\begin{multicols}{2}
			\begin{parts}
				% \part \includegraphics[width=0.5\linewidth]{mex_0013.png}
				
				% Perímetro: \\
				% \fillin[$P=(12\times 3) + (8\times 2)=36+16=52$][0in] \\
				
				% Área: \\
				% \fillin[$A=12^2+\dfrac{8^2}{2}=144+32=177$][0in]

				\part \includegraphics[width=.4\linewidth]{mex_0016.png}

				Perímetro: \\
				\fillin[$P=18\times 8 =144$][0in]

				Área: \\
				\fillin[$A=\dfrac{8\times 18 \times 9}{2}=648$][0in]

				\columnbreak%
				
				\part \includegraphics[width=.5\linewidth]{mex_0020.png}

				Perímetro: \\
				\fillin[$P=(2)61+(2)29=122+58=180$][0in]

				Área: \\
				\fillin[$A=61\times 29=1769$][0in]

			\end{parts}
		\end{multicols}
	}

	\questionboxed[3]{Encuentra el \textbf{perímetro} y el \textbf{área} de las siguientes figuras:

		\begin{multicols}{3}
			\begin{parts}\footnotesize
				% \part \includegraphics[width=.7\linewidth]{mex_0014.png}

				% Perímetro: \\
				% \fillin[$P=(3)24.51 + (2)12= 73.53+24=97.53$][0in]

				% Área: \\
				% \fillin[$A=24.51^2+\dfrac{12^2}{2}=600.74+72=672.74$][0in]

				% \part \includegraphics[width=.7\linewidth]{mex_0014.png} 

				% Perímetro: \\
				% \fillin[$P=(3)24.51 + (2)12= 73.53+24=97.53$][0in] 

				% Área: \\
				% \fillin[$A=24.51^2+\dfrac{12^2}{2}=600.74+72=672.74$][0in]

				% \part \includegraphics[width=.5\linewidth]{mex_0015.png}
				
				% Perímetro: \\
				% \fillin[$P=(18\times 3) + (11\times 2)=36+16=52$][0in] \\
				
				% Área: \\
				% \fillin[$A=18^2+\dfrac{11^2}{2}=324+60.5=384.5$][0in]

				\part \includegraphics[width=.5\linewidth]{mex_0019.png}
				
				Perímetro: \\
				\fillin[$P=(60\times 2) + (22\times 2)=164$][0in] \\
				
				Área: \\
				\fillin[$A=60\times 22=1320$][0in]

				% \part \includegraphics[width=.5\linewidth]{mex_0016.png} 

				% Perímetro: \\
				% \fillin[$P=18\times 8 =144$][0in] 

				% Área: \\
				% \fillin[$A=\dfrac{8\times 18 \times 9}{2}=648$][0in]

				% \part \includegraphics[width=.8\linewidth]{mex_0020.png} 

				% Perímetro: \\
				% \fillin[$P=(2)61+(2)29=122+58=180$][0in] 

				% Área: \\
				% \fillin[$A=61\times 29=1769$][0in]

				

				 \part \includegraphics[width=.5\linewidth]{mex_0017.png}
				
				 Perímetro: \\
				 \fillin[$P=(31\times 3) =93$][0in] \\
				
				 Área: \\
				 \fillin[$A=\frac{31\times 24}{2}=372$][0in]

				 \part \includegraphics[width=.5\linewidth]{mex_0018.png}
				
				 Perímetro: \\
				\fillin[$P=32\times 5 =160$][0in] \\
				
				Área: \\
				\fillin[$A=\dfrac{5\times 32 \times 12}{2}=960$][0in]
			\end{parts}
		\end{multicols}
	}